% chapter/appendix.tex
% Appendix – Quick Reference, Printable Tools, and Sealed Secrets
% Rewritten in the voice of Lerris Fenwood, with marginalia from the Grey Wanderer,
% Saikou Ira, and Dusana of the Raven Road.

\epigraph{“An appendix is not a graveyard for unused ideas. It is a toolkit for the desperate. Keep it close. Keep it honest. Keep it sealed when it needs to be.”}{—Lerris Fenwood, from the front matter of his own ledger}

What follows is not lore. It is **utility**. Tables you can glance at during play. Cards you can print and cut. A sealed section for the GM’s eyes only. And a final essay from me, Lerris Fenwood, on the counting of wounds – because I have counted more than I care to admit.

Use these pages when the players have taken a road you did not map. They will. They always do.

\section*{Appendix A – Quick Reference Tables}

\subsection*{Threat Level Summary}

\noindent (Repeated from Chapter 7 for ease of use.)

\begin{tabular}{|l|c|c|c|c|}
\hline
\textbf{Threat Level} & \textbf{Typical Use} & \textbf{Base Harm} & \textbf{Base Fatigue} & \textbf{Attribute Adj.} \\
\hline
0 & Non‑combat, flavour & 1–2 & 1 & –1 \\
1 & Standard minion & 3 & 2 & 0 \\
2 & Elite, lieutenant & 4–5 & 3 & +1 \\
3 & Mini‑boss & 6–7 & 4 & +2 \\
4+ & Major boss & 8+ & 5+ & +3 \\
\hline
\end{tabular}

\subsection*{Common Timer Sizes and Their Meaning}

\begin{tabular}{|l|l|}
\hline
\textbf{Segments} & \textbf{Typical Use} \\
\hline
4 & Short scene, minor threat, a single chase \\
6 & Standard encounter, a duel, a negotiation \\
8 & Extended fight, a dungeon level, a complex ritual \\
10 & Epic boss, a siege, the finale of a campaign \\
\hline
\end{tabular}

\noindent\textbf{Rule of Thumb}: A 4‑segment timer takes about 15–20 minutes of play. An 8‑segment timer takes 40–60 minutes. Adjust for table pace.

\subsection*{Harm and Fatigue at a Glance}

\begin{tabular}{|l|l|}
\hline
\textbf{Harm Level} & \textbf{Effect} \\
\hline
1 & –1 die on related actions \\
2 & –1 die on most actions until treated \\
3 & Incapacitated or dying \\
\hline
\end{tabular}

\noindent\textbf{Fatigue}: Each point worsens Position by one step (Dominant → Controlled → Desperate). If already Desperate, each Fatigue instead imposes –1 die. Full Fatigue track (Body rating) converts to Harm +1 and clears all Fatigue.

\subsection*{Armor Conversion (Quick)}

\begin{tabular}{|l|c|c|c|}
\hline
\textbf{Armor} & \textbf{Harm 1 →} & \textbf{Harm 2 →} & \textbf{Harm 3 →} \\
\hline
Light & Fatigue 1 & Harm 1 + Fatigue 1 & Harm 2 + Fatigue 1 \\
Medium & Fatigue 1 & Fatigue 1 & Harm 2 + Fatigue 1 \\
Heavy & Fatigue 1 & Fatigue 1 & Fatigue 2 \\
\hline
\end{tabular}

\grey{I have seen a player forget the armor conversion three times in one fight. The fourth time, I wrote it on an index card. The card saved his character’s life. Do not be clever. Be clear.}

\section*{Appendix B – Printable Adversary Card Template}

\noindent Print this page at 100\% scale. Cut along the hairlines. Fold and glue or laminate for durability.

\vspace{0.5cm}

\begin{center}
\fbox{\parbox{0.8\textwidth}{\centering
\textbf{Adversary Card – Front}\\
\vspace{0.3cm}
\texttt{[Name]}\\
\texttt{[Threat Level] [Creature Type]}\\
\texttt{Harm: ] [ Fatigue: ]}\\
\texttt{Traits: \_\_\_\_\_\_\_\_\_\_\_\_\_\_\_\_\_\_\_\_\_\_\_\_\_\_\_\_\_}\\
\texttt{Complications: \_\_\_\_\_\_\_\_\_\_\_\_\_\_\_\_\_\_\_\_\_\_\_}\\
\texttt{Clock: \_\_\_\_\_\_\_\_ (segments) – Tick on \_\_\_\_\_\_\_\_\_}
}}
\end{center}

\vspace{0.5cm}

\begin{center}
\fbox{\parbox{0.8\textwidth}{\centering
\textbf{Adversary Card – Back}\\
\vspace{0.3cm}
\texttt{Weakness: \_\_\_\_\_\_\_\_\_\_\_\_\_\_\_\_\_\_\_\_\_\_\_\_\_\_\_\_\_}\\
\texttt{Resolution (if not combat): \_\_\_\_\_\_\_\_\_\_\_\_\_\_\_\_\_\_\_}\\
\texttt{Drop: \_\_\_\_\_\_\_\_\_\_\_\_\_\_\_\_\_\_\_\_\_\_\_\_\_\_\_\_\_\_\_\_\_\_}\\
\texttt{Marginal note: \_\_\_\_\_\_\_\_\_\_\_\_\_\_\_\_\_\_\_\_\_\_\_\_\_\_\_\_}
}}
\end{center}

\noindent\textbf{How to use}: Fill in the blanks before the session or during play. Keep a stack of blank cards behind the screen. I use a soft pencil; the Grey Wanderer uses blood (do not copy him).

\saikou{I prefer a silver mirror. The card is a backup. The mirror never lies.}

\section*{Appendix C – Icon Legend}

\noindent Some creatures in this compendium use icon codes for quick reference. The following legend applies.

\begin{tabular}{|l|l|}
\hline
\textbf{Icon} & \textbf{Meaning} \\
\hline
[W] & Weakness – silver, fire, a true name, etc. \\
[C] & Timer – size and trigger \\
[L] & Lair – special environmental action \\
[P] & Phase – boss phase transition \\
[D] & Drop – treasure or boon \\
[WIT] & Requires a witness to resolve \\
[NIN] & Involves the ninth – counting, omission, taboo \\
\hline
\end{tabular}

\noindent Example: \emph{Slasher [W: silver] [C: Hunger 6 – ticks on kill]}. Means the slasher is vulnerable to silver and has a Hunger timer that ticks each time it kills.

\lerris{I invented this legend after a player asked me the same question four times in one fight. The legend did not help. The fight did not end well. But the legend was not at fault.}

\section*{Appendix D – Sealed Section (GM Eyes Only)}

\noindent The following pages are sealed. Open only when you need a secret adversary, a hidden weakness, or a final twist. The seal is not magical. It is a promise.

\subsection*{D.1 – The Unnamed Adversary: The Ledger Lich (Tier IV)}

\noindent\textbf{Description}: A curator bound by indices; hands of ink, breath of dust. It appears in archives, counting‑houses, and the little rooms where ledgers are kept. It does not fight. It **audits**.

\noindent\textbf{Stats}: Body 2, Wits 5, Spirit 5, Presence 4. Harm 6. Fatigue Spirit.

\noindent\textbf{Phase Timer }: [8] – Citation → Injunction → Foreclosure → Redaction

\noindent\textbf{Moves}: \emph{Name Forfeiture} (steal a declared tag until dawn), \emph{Writ of Seizure} (lock a Diamond or gear under spectral seal), \emph{Errata} (retroactively worsen one PC’s prior roll Outcome).

\noindent\textbf{Lair Actions}: \emph{Quiet Hours} (casting/rites generate +1 SB or fail softly), \emph{Reference Only} (movement through stacks costs an action unless proper tokens are presented), \emph{Cross‑reference Collapse} (bookshelves reconfigure; split the party).

\noindent\textbf{Counters}: True names carved in quartz, burning a citation (sacrifice an index – a map, a ledger, a vow) to cancel one Move.

\noindent\textbf{Drop}: A Key of Echoes (unlocks any door, but the door remembers) or a Tome of Convenient Knowledge (once per session, answer a single question about a creature’s weakness).

\subsection*{D.2 – The Unwritten Puzzle: The Ninth Name}

\noindent\textbf{Setup}: The party finds a sealed envelope with the words “Open when the ninth bell rings” written in a familiar hand (a dead mentor, a lost sibling). Inside is a single name.

\noindent\textbf{The Puzzle}: The name is not a name. It is a **count**. Each letter corresponds to a number (A=1, B=2, etc.). Sum the numbers. Divide by eight. The remainder is the number of the threshold where the envelope was sealed.

\noindent\textbf{The Secret}: The threshold is a door that has been bricked over in the Wardens’ archive. Behind it is a confession – the mentor’s final witness. The party must decide whether to read it aloud.

\noindent\textbf{Sealed Solution}: The name is \texttt{ALAYSE} (A=1, L=12, A=1, Y=25, S=19, E=5 → sum = 63. 63 ÷ 8 = 7 remainder 7. The threshold is the seventh door in the east corridor. The bricked‑over door. Behind it, the confession: \emph{“I lied about the ninth step. It is not a door. It is a witness. The witness is you.”}

\saikou{I helped seal this page. The ink is still warm. The secret is not a secret. It is a debt. Pay it when you are ready.}

\section*{Appendix E – On the Counting of Wounds (An Essay by Lerris Fenwood)}

\noindent I have killed seventeen creatures catalogued in this book. I have been wounded by twenty‑three. The wounds heal. The counting does not.

When I was young, my father taught me to count my breaths. “Eight,” he said. “Never nine. The ninth breath belongs to the Hollow.” I did not believe him. Then I saw a man take the ninth breath in an Aeler tunnel. The air turned cold. The man’s shadow stood up and walked away. The man did not follow. He simply sat down and waited. He is still waiting.

Adversaries are not bags of harm. They are **questions**. A ghoul asks: “What are you willing to eat?” A vampire asks: “What are you willing to lose?” A Oath‑Keeper asks: “What do you owe?” The dice do not answer these questions. The players do. The GM listens. The ledger is written in the space between.

I keep a second ledger for the wounds that did not heal. The first entry is a scar on my left hand – a ghoul’s claw. The second entry is a silence I cannot name – the debt I owe to the String‑Mother. The third entry is a door I did not open. I write this now so that you, reader, will know that the compendium is not complete. The compendium is a witness. The wounds are the testimony.

Count to eight. Write what you can prove. Leave the ninth line empty. The Hollow will fill it when it is ready.

\lerris{This essay is not advice. It is a confession. Use it as you will.}

\section*{Appendix F – Colophon}

\noindent This book was written on paper that had been a cargo manifest, a love letter, and a bandage. The ink is lamp‑black cut with salt water and the ash of a burnt confession. The binding is grey cord tied with nine knots. The ninth knot is not cut.

The Grey Wanderer contributed his usual cynicism and a recipe for tea that keeps the Hollow from your doorstep (mostly). Saikou Ira insisted on keeping a silver mirror beside the writing desk; it fogged exactly twice, and he will not say what he saw. Dusana sang a lark into the rafters; the bird stayed for three days and left a single feather, which I have pressed between these pages.

As for me, Lerris Fenwood, I have seen the faces of the Unmasked. I have read the true ledger of the Congregation. I have learned that the ninth name is never spoken because it belongs to the one who will speak it last. When that day comes, I hope this book serves as witness.

Until then, the road is south. The ink is fresh. The lamp is borrowed.

\grey{You forgot the bread. Always leave bread for the threshold.}
\saikou{And a coin for the ferryman, just in case.}
\dusana{And a story for the fire. The rest will take care of itself.}
\lerris{Paid in full.}

